\chapter{\ifproject%
\ifenglish Experimentation and Results\else การทดลองและผลลัพธ์\fi
\else%
\ifenglish System Evaluation\else การประเมินระบบ\fi
\fi}

บทนี้นำเสนอการทดลองเบื้องต้นที่ดำเนินการแล้ว 2 การทดลอง ได้แก่
(1) การส่ง Raw Logs ให้ LLM เชิงพาณิชย์วิเคราะห์โดยตรง เพื่อยืนยันปัญหาของโครงงาน และ
(2) การทดลองนำร่องเพื่อตรวจสอบว่า LLM ขนาดเล็กยังจำเป็นต้องใช้ RAG หรือไม่
จากนั้นนำเสนอแผนการทดลองหลัก ตัวชี้วัด ลำดับการดำเนินการ และประเด็นที่ยังเปิดอยู่

\section{การทดลองเบื้องต้น: ส่ง Raw Logs ให้ LLM วิเคราะห์โดยตรง}

\subsection{วิธีการทดลอง}
การทดลองนี้ดำเนินการก่อนการนำเสนอความก้าวหน้าครั้งที่ 1 โดยดึง Windows Sysmon Logs ของชุดข้อมูล T1087.001 จาก Splunk ผ่าน REST API
ได้เป็นไฟล์ JSON Lines จำนวน 11,090 Events (ข้อมูลต้นฉบับในรูปแบบ XML มีขนาดประมาณ 16.8 MB)
ซึ่งมีขนาดใหญ่กว่า Context Window ของ LLM ทั่วไป
แล้วอัปโหลดไฟล์ให้ LLM เชิงพาณิชย์ 3 โมเดลวิเคราะห์โดยตรง
ด้วยคำถามว่า ``นี่คือ Attack Technique อะไร เพราะอะไรถึงคิดแบบนั้น''
คำตอบที่ถูกต้อง (Ground Truth) คือ Account Discovery: Local Account (T1087.001)

\subsection{ผลการทดลอง}
\begin{table}[ht]
  \centering
  \caption{ผลการส่ง Raw Logs ของ T1087.001 ให้ LLM เชิงพาณิชย์วิเคราะห์โดยตรง}
  \label{tab:raw-llm}
  \begin{tabular}{|L{0.25\textwidth}|L{0.45\textwidth}|L{0.17\textwidth}|}
    \hline
    \textbf{โมเดล} & \textbf{คำตอบ} & \textbf{ผลเทียบเฉลย} \\ \hline
    GPT-5.6 Luna      & T1003 OS Credential Dumping & ผิด \\ \hline
    Gemini 3.1 Pro    & ไม่พบพฤติกรรมการโจมตีที่ชัดเจน (Logs เป็นเพียงการทำงานของ Splunk Agent) & ผิด \\ \hline
    Claude Sonnet 5 (High) & T1087 Account Discovery & ถูกเฉพาะระดับ Technique (Parent Match) \\ \hline
  \end{tabular}
\end{table}

จากตารางที่~\ref{tab:raw-llm} พบว่าแม้เป็นโมเดลขนาดใหญ่ แต่เมื่อได้รับ Raw Logs ที่มี Noise จำนวนมาก
บางโมเดลให้ความสำคัญกับ Events ที่ไม่เกี่ยวข้อง เช่น Gemini 3.1 Pro
ที่สรุปว่า Logs เป็นเพียงการทำงานของ Splunk Agent
และไม่มีโมเดลใดระบุ Sub-technique ได้ถูกต้อง
ผลนี้เป็นหลักฐานเบื้องต้นที่สนับสนุนว่าการสร้าง Investigation Context ที่คัดกรองแล้วมีความจำเป็น
ทั้งนี้ การทดลองนี้ทำผ่านหน้าเว็บของแต่ละบริการเพียงครั้งเดียว จึงควบคุมการตั้งค่าของโมเดลไม่ได้
และใช้เพียง 1 กรณี ผลจึงใช้เป็นตัวอย่างประกอบปัญหา ไม่ใช่การเปรียบเทียบความสามารถของโมเดล

\section{การทดลองนำร่อง: ความจำเป็นของ RAG สำหรับ LLM ขนาดเล็ก}

\subsection{วัตถุประสงค์และการตั้งค่า}
การทดลองนี้มีวัตถุประสงค์เพื่อตรวจสอบว่า เมื่อ Input ผ่าน Event Selection แล้ว
LLM ขนาดเล็ก (Qwen3.5:9B) ยังจำเป็นต้องใช้ RAG หรือไม่
และ Query รูปแบบใดเหมาะกับการค้นหาในฐานความรู้ ATT\&CK
โดยดำเนินการหลังการนำเสนอความก้าวหน้าครั้งที่ 1 และก่อนการนำเสนอครั้งที่ 2 (12--16 ตุลาคม 2569)

Input ของการทดลองคือผลลัพธ์จาก Event Selection แบบ Top-$k$ ของชุดข้อมูล T1087.001
ซึ่งลดจำนวน Events จาก 11,090 Events เหลือ 16 Events
(เทียบกับวิธี Multiple Related ที่ได้ 177 Events)
ประกอบด้วย Alert และ 16 Events ในรูปแบบ JSON ได้แก่ EID 1 จำนวน 3 Events คือ \texttt{cmd.exe},
\texttt{net.exe} (\texttt{net users}) และ \texttt{net1.exe} (\texttt{net1 users})
และ EID 10 จำนวน 13 Events
ในการทดลองนำร่อง Anchor ถูกกำหนดเองที่ Event ของคำสั่ง \texttt{net users}
ส่วนการทดลองหลักใช้ Alert ที่สร้างจาก ESCU ตามหัวข้อ~\ref{sec:alert}
โดยแต่ละ Event ใช้ประมาณ 100--130 Tokens (หัวข้อ Budget ในบทที่ 3)
การทดลองใช้ Embedding Model \texttt{qwen3-embedding:0.6b} และฐานความรู้ ATT\&CK v19.2 (697 Chunks)
ตามที่อธิบายในบทที่ 3
และตั้งค่า Temperature = 0 ปิดโหมด Thinking และ Context Window 16,384 Tokens เช่นเดียวกับการทดลองหลัก
ทั้งนี้ $k$ ของ RAG คือจำนวนเทคนิคที่ค้นคืนมา ซึ่งเป็นคนละความหมายกับ Top-$k$ ของ Event Selection
เงื่อนไขการทดลองแสดงในตารางที่~\ref{tab:rag-conditions}

\begin{table}[ht]
  \centering
  \caption{เงื่อนไขการทดลองความจำเป็นของ RAG}
  \label{tab:rag-conditions}
  \begin{tabular}{|L{0.25\textwidth}|c|L{0.52\textwidth}|}
    \hline
    \textbf{การทดลอง} & \textbf{RAG} & \textbf{คำอธิบาย} \\ \hline
    1. Direct-Log      & ไม่ใช้ & ส่ง Representation Log ให้ LLM โดยตรง \\ \hline
    2. Direct-Summary  & ไม่ใช้ & ให้ LLM แปลง Log เป็นคำอธิบายพฤติกรรมเป็นภาษาธรรมชาติ (ภาษาอังกฤษ) ก่อน แล้วจึงให้ระบุเทคนิค \\ \hline
    3.1 RAG-Log ($k=10$) & ใช้ & ใช้ Representation Log เป็น Query ค้นหา RAG ได้ Top-10 ให้ LLM \\ \hline
    3.2 RAG-Log ($k=5$)  & ใช้ & ใช้ Representation Log เป็น Query ค้นหา RAG ได้ Top-5 ให้ LLM \\ \hline
    4. RAG-Summary ($k=10$) & ใช้ & ใช้คำอธิบายพฤติกรรม (จากการทดลองที่ 2) เป็น Query ค้นหา RAG ได้ Top-10 ให้ LLM \\ \hline
  \end{tabular}
\end{table}

\subsection{ผลการทดลอง}
ผลการทดลองทั้ง 5 เงื่อนไข ซึ่งประกอบด้วยอันดับของเฉลยในผลการค้นคืนของ RAG และคำตอบของ LLM แสดงในตารางที่~\ref{tab:rag-results}

\begin{table}[ht]
  \centering
  \caption{ผลการทดลองความจำเป็นของ RAG (คำตอบที่ถูกต้องคือ T1087.001)}
  \label{tab:rag-results}
  \small
  \begin{tabular}{|l|c|l|c|}
    \hline
    \textbf{การทดลอง} & \textbf{อันดับเฉลยใน RAG} & \textbf{คำตอบของ LLM} & \textbf{ผล} \\ \hline
    1. Direct-Log        & --                & T1083 File and Directory Discovery & ผิด \\ \hline
    2. Direct-Summary    & --                & T1078 Valid Accounts                & ผิด \\ \hline
    3.1 RAG-Log ($k=10$) & อันดับ 6          & T1087.001                           & ถูก \\ \hline
    3.2 RAG-Log ($k=5$)  & ไม่อยู่ใน Top-5   & T1033 System Owner/User Discovery   & ผิด \\ \hline
    4. RAG-Summary ($k=10$) & อันดับ 1         & T1087.001                           & ถูก \\ \hline
  \end{tabular}
\end{table}

\subsection{การวิเคราะห์ผล}
จากตารางที่~\ref{tab:rag-results} พบว่า
\begin{enumerate}
  \item เมื่อไม่ใช้ RAG (การทดลองที่ 1 และ 2) LLM ขนาดเล็กตอบผิดทั้งสองกรณี
        แม้ Input จะผ่านการคัดเลือก Event มาแล้ว ชี้ว่าโมเดลขนาด 9B อาจมีความรู้ด้าน ATT\&CK
        ไม่เพียงพอที่จะระบุเทคนิคได้ตรง โดยตอบเป็นเทคนิคอื่นที่เกี่ยวข้องกับไฟล์หรือบัญชีผู้ใช้ (T1083 และ T1078)
  \item LLM ตอบถูกเฉพาะกรณีที่เฉลย (T1087.001) อยู่ในรายการ Top-$k$ จาก RAG
        แม้ Prompt จะระบุว่าคำตอบที่ถูกอาจไม่อยู่ในรายการ
        แสดงว่าในกรณีนี้ความถูกต้องของ LLM ขึ้นกับคุณภาพของการค้นคืน (Retrieval) เป็นหลัก
  \item การใช้คำอธิบายพฤติกรรมเป็น Query (RAG-Summary) ทำให้เฉลยขึ้นมาอยู่อันดับ 1
        ดีกว่าการใช้ JSON โดยตรง (อันดับ 6)
        สมมติฐานคือคำอธิบายภาษาธรรมชาติมีความหมายใกล้เคียงกับคำอธิบายเทคนิคในฐานความรู้มากกว่า JSON
        ซึ่งจะทดสอบกับกรณีจำนวนมากขึ้นในการทดลอง E3
  \item ค่า $k$ มีผลต่อผลลัพธ์ เมื่อลด $k$ จาก 10 เป็น 5 เฉลยซึ่งอยู่อันดับ 6 หลุดออกจากรายการ
        และ LLM ตอบผิด (การทดลองที่ 3.2)
        อย่างไรก็ตาม $k$ ที่มากขึ้นทำให้ Prompt ยาวขึ้นประมาณ 280 Tokens ต่อเทคนิค
\end{enumerate}

\subsection{ข้อจำกัดของการทดลอง}
การทดลองนี้ใช้ชุดข้อมูลเพียง 1 เทคนิค 1 กรณี และรันเงื่อนไขละ 1 ครั้ง
อีกทั้งคำสั่ง \texttt{net users} มีลักษณะชัดเจน จึงยังไม่สามารถสรุปผลโดยทั่วไปได้
ผลนี้ใช้เป็นแนวทางออกแบบการทดลอง E3 (หัวข้อ~\ref{sec:main-exp})
ซึ่งทดสอบกับทุกกรณีใน Test Set และกำหนดเกณฑ์การประเมินไว้ล่วงหน้า

\section{แผนการทดลองหลัก}
\label{sec:main-exp}

การทดลองหลักใช้ Development Set และ Test Set ตามหัวข้อ~\ref{sec:dataset}
โดยปรับพารามิเตอร์และ Prompt บน Development Set แล้วล็อกค่า (Freeze)
ก่อนรันบน Test Set เพียงครั้งเดียว

\subsection{Baseline และหลักการเปลี่ยนทีละปัจจัย}
Baseline ของการทดลองคือ Event Selection แบบ Top-$k$ ที่งบประมาณ $B^{*}$
นำเสนอในรูปแบบ JSON และส่งให้ LLM โดยตรงโดยไม่ใช้ RAG (Direct-Log)
โดย $B^{*}$ คือค่า $B \in \{500, 1{,}000, 2{,}000, 4{,}000\}$ Tokens ที่ได้ Exact Match Accuracy สูงสุดบน Development Set (หากเท่ากันเลือกค่า $B$ ที่น้อยกว่า)
ค่า $B$ สูงสุดกำหนดที่ 4,000 Tokens เพราะ Context Window ขนาด 16,384 Tokens
ต้องเผื่อพื้นที่ให้ Prompt ผลลัพธ์จาก RAG ($k=10$ ประมาณ $10 \times 280$ Tokens) และคำตอบของโมเดล
และการเพิ่มค่าทีละสองเท่าช่วยให้เห็นจุดที่ความถูกต้องหยุดเพิ่มขึ้น

การทดลองใช้หลักเปลี่ยนทีละปัจจัยจาก Baseline (One-factor-at-a-time)
เนื่องจากการทดลองทุกการผสม (Event Selection 5 วิธี $\times$ Representation 5 รูปแบบ $\times$ เงื่อนไขของ E3 5 แบบ $= 125$ รูปแบบ $\times$ 42 กรณี $= 5{,}250$ ครั้ง)
ใช้เวลามากเกินกว่าที่โครงงานทำได้ ข้อจำกัดคือไม่เห็นปฏิสัมพันธ์ระหว่างปัจจัย

\subsection{R0: Raw Logs}
เพื่อตอบวัตถุประสงค์ข้อที่ 3 ว่า Investigation Context ทำให้ LLM ระบุเทคนิคได้ถูกต้องกว่าการใช้ Raw Logs หรือไม่
R0 ส่ง Alert พร้อม Events ดิบทั้งหมดในช่วงเวลา $\pm T$ รอบ Alert ในรูปแบบต้นฉบับ
โดยไม่คัดเลือก ไม่ Normalize และไม่ตัด Field ออก
หากข้อมูลยาวเกิน Context Window จะลดค่า $T$ จนพอดีและบันทึกค่า $T$ ที่ใช้จริง

\subsection{E1: Event Selection}
เปรียบเทียบวิธี Event Selection ทั้ง 5 วิธี โดยใช้รูปแบบ JSON และไม่ใช้ RAG (ตารางที่~\ref{tab:exp-e1})
วิธีอื่นนอกจาก Top-$k$ ไม่จำกัดด้วยงบประมาณ $B$ แต่บันทึกจำนวน Token ทุกครั้ง

\begin{table}[ht]
  \centering
  \caption{รูปแบบการทดลอง E1: Event Selection}
  \label{tab:exp-e1}
  \small
  \begin{tabular}{|l|l|L{0.5\textwidth}|}
    \hline
    \textbf{รหัส} & \textbf{วิธี} & \textbf{สิ่งที่ส่งให้ LLM} \\ \hline
    E1-1 & Alert Only & เฉพาะ Field ของ Alert \\ \hline
    E1-2 & Single Event & ทุก Field ของ Event ที่ Alert ชี้ (1 Event) \\ \hline
    E1-3 & Time Window & Events ที่ Normalize แล้วทั้งหมดในช่วง $\pm T$ \\ \hline
    E1-4 & Multiple Related & Events ในสายความสัมพันธ์ ProcessGuid (บรรพบุรุษและลูกหลาน) ไม่ให้คะแนนและไม่จำกัดงบประมาณ \\ \hline
    E1-5 ถึง E1-8 & Top-$k$ & Graph $\rightarrow$ Traverse $\rightarrow$ Score $\rightarrow$ Budget ที่ $B$ = 500, 1,000, 2,000 และ 4,000 \\ \hline
  \end{tabular}
\end{table}

\subsection{E2: Representation}
ใช้ Top-$k$ ที่ $B^{*}$ และไม่ใช้ RAG แล้วเปลี่ยนเฉพาะรูปแบบการนำเสนอข้อมูล (ตารางที่~\ref{tab:exp-e2})
การคัดเลือก Events ใช้งบประมาณที่นับเป็น Token ของ JSON แล้วแปลงชุด Events เดียวกันเป็นรูปแบบอื่น
ทุกรูปแบบจึงมีข้อมูลชุดเดียวกัน ข้อจำกัดคือรูปแบบ Timeline ตัด Field บางส่วนออก

\begin{table}[ht]
  \centering
  \caption{รูปแบบการทดลอง E2: Representation}
  \label{tab:exp-e2}
  \small
  \begin{tabular}{|l|l|c|}
    \hline
    \textbf{รหัส} & \textbf{รูปแบบ} & \textbf{Tokens ต่อ Event (ตารางที่~\ref{tab:representation})} \\ \hline
    Baseline & JSON & $\approx$ 138 \\ \hline
    E2-1 & XML & $\approx$ 195 \\ \hline
    E2-2 & TOON & $\approx$ 119 \\ \hline
    E2-3 & Key=Value & $\approx$ 120 \\ \hline
    E2-4 & Timeline & $\approx$ 35 \\ \hline
  \end{tabular}
\end{table}

\subsection{E3: RAG}
ใช้ Top-$k$ ที่ $B^{*}$ และรูปแบบ JSON แล้วเปลี่ยนวิธีการใช้ RAG (ตารางที่~\ref{tab:exp-e3})
ซึ่งขยายจากการทดลองนำร่องในหัวข้อก่อนหน้าให้ครอบคลุมทุกกรณีใน Test Set
ย่อหน้าสรุปพฤติกรรม (Summary) สร้างครั้งเดียวต่อกรณีและใช้ร่วมกันใน E3-1 และ E3-4
โดยนับจำนวน Token และเวลาของการสร้างย่อหน้าสรุปรวมไว้ด้วย
เงื่อนไข RAG-Summary ทดลองเฉพาะ $k=10$ เพื่อจำกัดจำนวนรูปแบบการทดลอง
ทั้งนี้ผลการค้นคืนชุดเดียวกันยังใช้วัด hit@5 ของ Query แบบ Summary ได้โดยไม่ต้องเรียก LLM เพิ่ม

\begin{table}[ht]
  \centering
  \caption{รูปแบบการทดลอง E3: RAG}
  \label{tab:exp-e3}
  \small
  \begin{tabular}{|l|l|L{0.22\textwidth}|L{0.28\textwidth}|}
    \hline
    \textbf{รหัส} & \textbf{ชื่อ} & \textbf{Query ของ RAG} & \textbf{สิ่งที่ LLM เห็น} \\ \hline
    Baseline & Direct-Log & -- & Logs \\ \hline
    E3-1 & Direct-Summary & -- & ย่อหน้าสรุป 80--150 คำ (ไม่มีชื่อหรือรหัสเทคนิค) \\ \hline
    E3-2 & RAG-Log ($k=10$) & Logs JSON ทั้งหมด & Logs + 10 เทคนิคที่ค้นคืนได้ \\ \hline
    E3-3 & RAG-Log ($k=5$) & Logs JSON ทั้งหมด & Logs + 5 เทคนิคที่ค้นคืนได้ \\ \hline
    E3-4 & RAG-Summary ($k=10$) & ย่อหน้าสรุป & Logs + 10 เทคนิคที่ค้นคืนได้ \\ \hline
  \end{tabular}
\end{table}

\subsection{จำนวนการทดลองและค่าที่ควบคุม}
การทดลองบน Test Set มีทั้งหมด 17 รูปแบบ (ตารางที่~\ref{tab:exp-runs})
ค่าที่คงที่ตลอดการทดลอง ได้แก่ LLM Qwen3.5:9B ผ่าน Ollama, Temperature = 0, ปิดโหมด Thinking,
Context Window 16,384 Tokens, Embedding Model \texttt{qwen3-embedding:0.6b},
ฐานความรู้ ATT\&CK v19.2 และ Prompt แม่แบบเดียวกัน
(เงื่อนไขที่ใช้ RAG เพิ่มเฉพาะส่วนรายการ Candidate)
ส่วนค่าที่ปรับบน Development Set แล้วล็อกไว้ ได้แก่ $\gamma$, $\tau$, $w[\mathrm{EventID}]$,
จำนวน Hop สูงสุด, $B^{*}$, ค่า $\pm T$ ของ Time Window และ R0 รวมถึง Prompt

\begin{table}[ht]
  \centering
  \caption{จำนวนรูปแบบการทดลองและจำนวนครั้งที่รันบน Test Set}
  \label{tab:exp-runs}
  \small
  \begin{tabular}{|l|c|c|}
    \hline
    \textbf{ชุดการทดลอง} & \textbf{จำนวนรูปแบบ} & \textbf{จำนวนครั้ง (42 กรณี)} \\ \hline
    R0 & 1 & 42 \\ \hline
    E1 (รวม Baseline) & 8 & 336 \\ \hline
    E2 & 4 & 168 \\ \hline
    E3 & 4 & 168 \\ \hline
    \textbf{รวม} & \textbf{17} & \textbf{714 (+ 42 ครั้งสำหรับสร้างย่อหน้าสรุป)} \\ \hline
  \end{tabular}
\end{table}

\subsection{ตัวชี้วัดและการวิเคราะห์ทางสถิติ}
ตัวชี้วัดที่ใช้ตามหัวข้อ~\ref{sec:metrics} จับคู่กับวัตถุประสงค์ของโครงงานดังตารางที่~\ref{tab:exp-metrics}
ค่าความถูกต้องจะรายงานพร้อมช่วงความเชื่อมั่น 95\% ด้วยวิธี Bootstrap
และเปรียบเทียบแต่ละรูปแบบกับ Baseline ด้วย McNemar's Test จำนวน 16 คู่ โดยปรับค่า $p$ ด้วยวิธีของ Holm
นอกจากนี้จะรายงานผลแยกตามเทคนิค แยกกรณีที่ใช้ Anchor สำรองออกจากกรณีที่มี Alert จริง
และแยกตามกรณี A/B/C ของ Lineage (หัวข้อ~\ref{sec:traverse})

คำตอบของ LLM อยู่ในรูปแบบ JSON ตามที่กำหนดใน Prompt ระบบจะดึงรหัสเทคนิคจากคำตอบมาเทียบกับเฉลย
หาก LLM ตอบผิดรูปแบบหรือไม่ระบุรหัสเทคนิค จะนับเป็นคำตอบที่ผิด และรายงานจำนวนกรณีดังกล่าวแยกไว้
เนื่องจากป้ายกำกับของทั้งสองแหล่งอาจสร้างขึ้นก่อน ATT\&CK v19.2 จะตรวจสอบและแปลงรหัสเทคนิค
ที่ถูกเปลี่ยนหรือยกเลิก (Revoked/Deprecated) ให้เป็นรหัสใน v19.2 ก่อนล็อกค่าการทดลอง

\begin{table}[ht]
  \centering
  \caption{ตัวชี้วัดที่ใช้ตอบวัตถุประสงค์ของโครงงาน}
  \label{tab:exp-metrics}
  \small
  \begin{tabular}{|L{0.3\textwidth}|L{0.38\textwidth}|l|}
    \hline
    \textbf{วัตถุประสงค์} & \textbf{ตัวชี้วัด} & \textbf{วัดบน} \\ \hline
    ข้อ 2: เปรียบเทียบ Event Selection และ Representation & Accuracy และ Input Tokens ของ E1 และ E2 & Test Set \\ \hline
    ข้อ 3: ความถูกต้องในการระบุเทคนิค & Exact Match และ Parent Match Accuracy & Test Set \\ \hline
    ข้อ 3: ประสิทธิภาพของ RAG (E3) & อันดับของเฉลย, hit@5 และ hit@10 & Test Set \\ \hline
    ข้อ 4: การลดปริมาณข้อมูลและต้นทุน & Input Tokens, Latency และสัดส่วน Events ที่ถูกคัดเลือก & Test Set \\ \hline
    ข้อ 4: การคงหลักฐานที่จำเป็น & Evidence Recall & Development Set \\ \hline
    ข้อมูลประกอบปัญหา & จำนวน Alert ต่อไฟล์ และจำนวน TP/FP & ทั้งสองชุด \\ \hline
  \end{tabular}
\end{table}

\section{ลำดับการดำเนินการทดลอง}
การทดลองแบ่งเป็น 5 ระยะ ดังตารางที่~\ref{tab:exp-phases}
โดย Test Set ผ่านเพียงการเตรียมข้อมูลในระยะที่ 1
การปรับพารามิเตอร์ในระยะที่ 2--3 ใช้เฉพาะ Development Set
และ Test Set ถูกใช้ประเมินผลเพียงครั้งเดียวในระยะที่ 4

\begin{table}[ht]
  \centering
  \caption{ลำดับการดำเนินการทดลอง}
  \label{tab:exp-phases}
  \small
  \begin{tabular}{|c|L{0.55\textwidth}|l|}
    \hline
    \textbf{ระยะ} & \textbf{งาน} & \textbf{ข้อมูล} \\ \hline
    0 & การทดลองนำร่อง (เสร็จสิ้น) & T1087.001 \\ \hline
    1 & Freeze, Normalize, Screening, สร้าง Alert และกรณีทดสอบ & Dev + Test \\ \hline
    2 & สร้าง Ground Truth ระดับ Event อัตโนมัติจาก Process Tree ของ Atomic Red Team แล้วสุ่มตรวจ & Dev \\ \hline
    3 & ปรับพารามิเตอร์แล้วล็อกค่า & Dev \\ \hline
    4 & รันทั้ง 17 รูปแบบเพียงครั้งเดียว & Test \\ \hline
  \end{tabular}
\end{table}

\section{ประเด็นที่ยังเปิดและข้อจำกัด}
\begin{itemize}
  \item \textbf{การเลือก Alert เมื่อมีหลายรายการ} -- รอบนี้ใช้ 1 Alert ต่อไฟล์ การเลือกหรือจัดกลุ่ม Alert จะศึกษาต่อ
  \item \textbf{ป้ายกำกับที่น่าสงสัยใน OTRF} -- เช่น ชุด Rubeus ASKTGT ที่ระบุเป็น T1003.003
        และชุด MSBuild ที่ระบุเป็น T1047 (หัวข้อ~\ref{sec:dev-test}) ต้องทบทวนก่อนล็อกค่าการทดลอง
  \item \textbf{ตัวแปลง OTRF JSON เป็น Sysmon XML} -- ยังต้องพัฒนา และเป็นความเสี่ยงหลักของการสร้าง Alert ฝั่ง Test Set
  \item \textbf{จำนวนกรณีทดสอบอาจลดลง} -- ชุดข้อมูล OTRF ตรวจแล้วว่ามี Sysmon Event
        แต่ยังต้องเทียบกับคำสั่งใน Adversary View
  \item \textbf{ขนาดของ Test Set} -- 42 กรณีทำให้ช่วงความเชื่อมั่นค่อนข้างกว้าง
        ความแตกต่างระหว่างรูปแบบการทดลองที่มีขนาดเล็กอาจไม่มีนัยสำคัญทางสถิติ
  \item \textbf{T1059.001 (PowerShell)} -- Prompt ปัจจุบันให้ LLM ถือว่า PowerShell หรือ Command Shell
        เป็นเพียงช่องทาง (Carrier) ของเทคนิคอื่น เมื่อเฉลยคือ T1059.001 โมเดลจึงอาจตอบเป็นเทคนิคอื่นตาม Prompt
        ผลของเทคนิคนี้จะรายงานแยกไว้
  \item \textbf{Ground Truth อยู่ในระดับชุดข้อมูล} -- ชุดข้อมูลระบุเพียงว่าไฟล์เป็นเทคนิคใด
        ไม่ได้ระบุว่า Event ใดเป็นการโจมตี จึงวัด Evidence Recall ได้เฉพาะ Development Set
  \item \textbf{การรั่วไหลของเฉลย} -- ในการทดลองนำร่องพบว่ารหัส Alert มีรหัสเทคนิคติดมาด้วย
        และ Alert จาก ESCU มีชื่อ Rule และป้ายกำกับ ATT\&CK
        จึงต้องทำ Leak Check กับทุก Prompt ก่อนส่งให้ LLM
  \item \textbf{ต้นแบบ Web Dashboard} -- ไม่อยู่ในการประเมินเชิงปริมาณของการทดลองนี้
  \item \textbf{นอกขอบเขต (งานในอนาคต)} -- การจำแนก True Positive/False Positive,
        การประเมินบน Linux Auditd, ชุดข้อมูลแบบ Compound หรือ APT
        และชุดข้อมูลของ Splunk ที่ไม่ได้สร้างด้วย Atomic Red Team
\end{itemize}
